\BeleskeID{07-s016}
% element: 07-s016:d01
% element: 07-s016:e02
Na početku velikog projekta VHDL može opisati performanse i zahteve interfejsa pojedinih komponenti, dok je projektovanje još na visokom nivou. To je posebno korisno kada radi više članova tima. Projektant sistema pristupom odozgo nadole definiše interfejs svake komponente i zahteve prihvatanja u obliku testbench-a visokog nivoa. Deklaracija entiteta i specifikacija performansi potom se daju drugim članovima tima da opis dovrše ili usavrše.

% element: 07-s016:e03
Pri unosu dizajna detalji sistema ulaze u računarski sistem za projektovanje. Mogu se uneti kao šeme na nivou ploče ili funkcionalne šeme, ili kao VHDL opisi. Ako se koristi sinteza, VHDL deo treba pisati stilom pogodnim za sintezu. Ova faza može koristiti i alate koji nisu VHDL; kompletan sistem često nastaje kombinovanjem VHDL opisa i drugih prikaza, uključujući šeme.
